\chapter{Advanced colspec}\label{app:colspec}

In this tutorial, we saw that \snippet{colspec} could contain information on the colunms and vertical borders of your table. But, just like with native Latex tables, there is more to it. In most cases however, tabularray has a header option that is preferrable.

\section{Column separation}

The directive \snippet{@\{\}} allows to specify what Latex must insert between two columns. For instance, if you want two columns to be separated by \snippet{2cm}, you can write:

\begin{Code}[show-output]
\begin{tblr}{
	colspec={l\highLight{@{\hspace{2cm}}}cr},
}
	a & b & c
\end{tblr}
\end{Code}

It is often better to rather use the \snippet{colsep} (margin on both sides), \snippet{leftsep}, \snippet{rightsep} options from tabularray, as they have the advantage of not being completely counter-intuitive: the \snippet{@\{\}} is incompatible with a vertical border, and the extra space does not belong to any column, therefore cannot be styled:

\begin{Code}[show-output]
\begin{tblr}{
	colspec={l@{\hspace{2cm}}c@{\hspace{2cm}}r},
	column{1} = {green!20},
	column{2} = {red!20},
	column{3} = {blue!20},
	vline{2}, % will break spacing between the first two columns
}
	a & b & c
\end{tblr}	
\end{Code}

\begin{Code}[show-output]
\begin{tblr}{
	colspec={lcr},
	column{1} = {rightsep=1cm, green!20},
	column{2} = {colsep=1cm, red!20},
	column{3} = {leftsep=1cm, blue!20},
	vline{2} = {},  % does not break anything
}
	a & b & c
\end{tblr}
\end{Code}

\section{Text at the beginning and the end of a cell}

Let us mention the operators \snippet{>\{\}} and \snippet{<\{\}}, who respectively add some text at the beginning and at the end of each cell of the column. Rather than long explanations, here is an example: 

\begin{Code}[show-output]
\begin{tblr}{\highLight{>{Candidate }}ccc\highLight{<{s}}}
      & Name   & Time in~\\ \hline
	1 & Claude & 10.5 \\
	2 & Mary   & 5.68 \\
	3 & Jorgen & 12.8 \\
\end{tblr}
\end{Code}

Here again, tabularray has the options \snippet{preto} and \snippet{appto} that are preferrable and do same thing in most cases. If you want to add commands with the operators  \snippet{>} and \snippet{<}, you will maybe have to use the option \snippet{cmd} instead. I will let you refer to the documentation for this one.

\section{Several times the same column}

When our table has numerous columns, writing dozens of times the same instruction in the colspec can quickly become repetitive. Luckily, you can replace this by, for instance, \snippet{*\{3\}\{c|Q[l,2cm]\}}, which is the same thing as \snippet{c|Q[l,2cm]c|Q[l,2cm]c|Q[l,2cm]}.
Here is an example, with as a bonus a use of the option \snippet{cmd} that we talked about just above.

\begin{Code}
\documentclass{standalone}
\usepackage{tabularray, xcolor, ninecolors}
\usepackage{adjustbox}  % for \adjustbox that rotates a textbox
\usepackage{bbding}  % for \Checkmark and \XSolidBrush
\newcommand{\yes}{\textcolor{green!70!black}{\Checkmark}}
\newcommand{\no}{\textcolor{red!50}{\XSolidBrush}}

\begin{document}
\begin{tblr}{
	colspec={Q[2cm,r,m]\highLight{*{27}{c}}},  % One 2cm column, and 27 centered ones
	columns = {font={\small},colsep=3.5pt},
	row{1} = {cmd=\adjustbox{angle=45,lap=\width-(1em)}},
	% the cmd option applies the command adjustbox (rotation) to every cell in the first row
}
	& Austria & Belgium & Bulgaria & Croatia & Cyprus & Czechia & Denmark & Estonia & Finland & France & Germany & Greece & Hungary & Ireland & Italy & Latvia & Lithuania & Luxembourg & Malta & Netherlands & Poland & Portugal & Romania & Slovakia & Slovenia & Spain & Sweden
	Eurozone & \yes & \yes & \yes & \yes & \yes & \yes & \no & \yes & \yes & \yes & \yes & \yes & \no & \yes & \yes & \yes & \yes & \yes & \yes & \yes & \no & \yes & \no & \yes & \yes & \no & \no \\
	Schengen Area & \yes & \yes & \yes & \yes & \no & \yes & \yes & \yes & \yes & \yes & \yes & \yes & \yes & \no & \yes & \yes & \yes & \yes & \yes & \yes & \yes & \yes & \yes & \yes & \yes & \yes & \yes & \\ 
	European Economic Area & \yes & \yes & \yes & \yes & \yes & \yes & \yes & \yes & \yes & \yes & \yes & \yes & \yes & \yes & \yes & \yes & \yes & \yes & \yes & \yes & \yes & \yes & \yes & \yes & \yes & \yes & \yes & \\ 
\end{tblr}
\hspace{2em} % a bit of extra space to avoid cutting the end of our table
\end{document}
\end{Code}

Will give:

\begin{center}
\begin{standalonebox}
\begin{tblr}{
	colspec={Q[2cm,r,m]*{27}{c}},  % One 2cm column, and 27 centered ones
	columns = {font={\small},colsep=3.5pt},
	row{1} = {cmd=\adjustbox{angle=45,lap=\width-(1em)}},
	% the cmd option applies the command adjustbox (rotation) to every cell in the first row
}
	& Austria & Belgium & Bulgaria & Croatia & Cyprus & Czechia & Denmark & Estonia & Finland & France & Germany & Greece & Hungary & Ireland & Italy & Latvia & Lithuania & Luxembourg & Malta & Netherlands & Poland & Portugal & Romania & Slovakia & Slovenia & Spain & Sweden \\
	Eurozone & \yes & \yes & \yes & \yes & \yes & \yes & \no & \yes & \yes & \yes & \yes & \yes & \no & \yes & \yes & \yes & \yes & \yes & \yes & \yes & \no & \yes & \no & \yes & \yes & \no & \no \\
	Schengen Area & \yes & \yes & \yes & \yes & \no & \yes & \yes & \yes & \yes & \yes & \yes & \yes & \yes & \no & \yes & \yes & \yes & \yes & \yes & \yes & \yes & \yes & \yes & \yes & \yes & \yes & \yes & \\ 
	European Economic Area & \yes & \yes & \yes & \yes & \yes & \yes & \yes & \yes & \yes & \yes & \yes & \yes & \yes & \yes & \yes & \yes & \yes & \yes & \yes & \yes & \yes & \yes & \yes & \yes & \yes & \yes & \yes & \\ 
\end{tblr}
\end{standalonebox}
\end{center}

\section{No colspec}

I have to confess something: the colspec option is not mandatory. As a matter of fact, tabularray computes the table structure from the ampersands \snippet{\&} in each row. Moreover, in case of conflict between the colspect and the table body, the table body wins (remember the priority rules!).

\begin{Code}
\begin{tblr}\highLight{{}}  % also works
	Hello & tabularray & world! \\
	cell & other cell & last cell \\
\end{tblr}
\end{Code}

The colspec option is essentially useful to give the reader a quick indication about the table structure, and is also useful for applying basic styling to each column. I strongly encourage you to always write the colspec, but know that if there are not enough columns in the compiled table, it's probably because you forgot a few \snippet{\&} somewhere!

Finally, the colspec is mandatory with native Latex tables, so if one day you have to migrate your tables back to native, you will know what to do.